\paragraph{Phase Kickback with CNOT}\label{ex:phase-kickback-with-cnot}

Until now, everything has been abstract:
\begin{itemize}
    \item Arbitrary unitary $U$;
    \item Arbitrary eigenstate \ket{v};
    \item Arbitrary eigenvalue $e^{i\phi}$.
\end{itemize}
Now we'll replace $U$ with one of the most familiar quantum gates: Pauli-$X$. This example is particularly important because \textbf{CNOT is nothing more than a controlled-$X$ gate}, and CNOT appears everywhere in quantum algorithms, including Deutsch's algorithm (that we'll see in the next sections).

\highspace
\begin{flushleft}
    \textcolor{Green3}{\faIcon{question-circle} \textbf{Remember what CNOT really is}}
\end{flushleft}
The CNOT gate performs the operation:
\begin{equation*}
    CX = \begin{pmatrix}
        I & 0 \\
        0 & X
    \end{pmatrix} =
    \begin{pmatrix}
        1 & 0 & 0 & 0 \\
        0 & 1 & 0 & 0 \\
        0 & 0 & 0 & 1 \\
        0 & 0 & 1 & 0
    \end{pmatrix}
\end{equation*}
This means:
\begin{itemize}
    \item If the control is \ket{0}, do nothing to the target;
    \item If the control is \ket{1}, apply $X$ to the target.
\end{itemize}
So CNOT is simply \textbf{a controlled Pauli-$X$ gate}.

\highspace
Now, phase kickback only works when the target is an eigenstate of the unitary. So we first ask: \emph{\textbf{what are the eigenstates of $X$?}} We know that $X$ is a Pauli matrix:
\begin{equation*}
    X = \begin{pmatrix}
        0 & 1 \\
        1 & 0
    \end{pmatrix}
\end{equation*}
To find the eigenstates:
\begin{itemize}
    \item Write a generic qubit:
    \begin{equation*}
        \ket{v} = \begin{pmatrix}
            a \\ b
        \end{pmatrix}
        \quad
        a,b \in \mathbb{C}
    \end{equation*}
    \item Apply $X$ to \ket{v}:
    \begin{equation*}
        X\ket{v} = \begin{pmatrix}
            0 & 1 \\
            1 & 0
        \end{pmatrix}
        \begin{pmatrix}
            a \\ b
        \end{pmatrix} =
        \begin{pmatrix}
            b \\ a
        \end{pmatrix}
    \end{equation*}
    \item Apply the eigenvalue definition:
    \begin{equation*}
        X\ket{v} = \lambda\ket{v} \quad \Rightarrow \quad
        \begin{pmatrix}
            b \\ a
        \end{pmatrix} =
        \lambda
        \begin{pmatrix}
            a \\ b
        \end{pmatrix}
        \quad \Rightarrow \quad
        \begin{cases}
            b = \lambda a \\
            a = \lambda b
        \end{cases}
    \end{equation*}
    \item Find the eigenvalues by substituting $b$ in the second equation:
    \begin{equation*}
        a = \lambda b = \lambda (\lambda a) = \lambda^2 a
    \end{equation*}
    Assuming $a \neq 0$, we can divide both sides by $a$:
    \begin{equation*}
        1 = \lambda^2 \quad \Rightarrow \quad \lambda = \pm 1
    \end{equation*}
    \item Find the eigenvectors by substituting the eigenvalues in the first equation:
    \begin{itemize}
        \item For $\lambda = 1$:
        \begin{equation*}
            b = \lambda a = 1 \cdot a = a \quad \Rightarrow \quad \ket{v_1} = \begin{pmatrix}
                a \\ a
            \end{pmatrix} = a\begin{pmatrix}
                1 \\ 1
            \end{pmatrix}
        \end{equation*}
        Eigenvectors can be multiplied by any nonzero costant and still represent the same direction, so we normalize it:
        \begin{equation*}
            \sqrt{1^2 + 1^2} = \sqrt{2} \quad \Rightarrow \quad \ket{v_1} = \frac{1}{\sqrt{2}}\begin{pmatrix}
                1 \\ 1
            \end{pmatrix} = \ket{+}
        \end{equation*}
        \item For $\lambda = -1$:
        \begin{equation*}
            b = \lambda a = -1 \cdot a = -a \quad \Rightarrow \quad \ket{v_2} = \begin{pmatrix}
                a \\ -a
            \end{pmatrix} = a\begin{pmatrix}
                1 \\ -1
            \end{pmatrix}
        \end{equation*}
        Normalizing it:
        \begin{equation*}
            \sqrt{1^2 + (-1)^2} = \sqrt{2} \quad \Rightarrow \quad \ket{v_2} = \frac{1}{\sqrt{2}}\begin{pmatrix}
                1 \\ -1
            \end{pmatrix} = \ket{-}
        \end{equation*}
    \end{itemize}
\end{itemize}
Therefore, the eigenstates of $X$ are \ket{+} and \ket{-}, with eigenvalues $+1$ and $-1$, respectively. The phase for each eigenvalue is:
\begin{itemize}
    \item For $\lambda = +1$: $e^{i\phi} = +1 \quad \Rightarrow \quad \phi = 0$;
    \item For $\lambda = -1$: $e^{i\phi} = -1 \quad \Rightarrow \quad e^{i\phi} = \cos (\phi) + i \sin (\phi) = -1 \quad \Rightarrow \quad \phi = \pi$.
\end{itemize}
With these phases, we can now see how phase kickback works with CNOT. We expect that if the target is in \ket{+}, the control will remain unchanged (i.e., no phase kickback, $\phi = 0$), while if the target is in \ket{-}, the control will acquire a phase of $\pi$. Mathematically, we can write:
\begin{equation*}
    CX\ket{c}\ket{+} = e^{i \cdot 0}\ket{c}\ket{+} = 1 \cdot \ket{c}\ket{+} = \ket{c}\ket{+} \quad \text{(no phase kickback)}
\end{equation*}
\begin{equation*}
    CX\ket{c}\ket{-} = e^{i\pi}\ket{c}\ket{-} = -\ket{c}\ket{-} \quad \text{(phase kickback of $\pi$)}
\end{equation*}

\highspace
\begin{examplebox}[: Target in \ket{+}]
    A $CX$ gate with the target in \ket{+} will kick back a phase of $0$ to the control qubit. This means that the control qubit will remain unchanged after the operation. Let's see if this is true:
    \begin{itemize}
        \item Control in \ket{+} to allow for superposition of control states:
        \begin{equation*}
            \ket{+} = \frac{\ket{0} + \ket{1}}{\sqrt{2}}
        \end{equation*}
        \item Target in \ket{+} to be an eigenstate of $X$ with eigenvalue $+1$:
        \begin{equation*}
            \ket{+} = \frac{\ket{0} + \ket{1}}{\sqrt{2}}
        \end{equation*}
    \end{itemize}
    The initial state is:
    \begin{equation*}
        \ket{\psi} = \ket{+}\ket{+} = \frac{\ket{0} + \ket{1}}{\sqrt{2}} \otimes \ket{+}
    \end{equation*}
    Applying the $CX$ gate (CNOT):
    \begin{center}
        \begin{quantikz}
            \lstick{\ket{+}} & \ctrl{1} & \rstick{?} \\
            \lstick{\ket{+}} & \gate{X} & \rstick{?}
        \end{quantikz}
    \end{center}
    \begin{equation*}
        CX\ket{\psi} = CX\left(\frac{\ket{0} + \ket{1}}{\sqrt{2}} \otimes \ket{+}\right) = \frac{\ket{0} \ket{+} + \ket{1} X \ket{+}}{\sqrt{2}}
    \end{equation*}
    Since $X\ket{+} = \ket{+}$, we have:
    \begin{equation*}
        CX\ket{\psi} = \frac{\ket{0} \ket{+} + \ket{1} \ket{+}}{\sqrt{2}} = \frac{\ket{0} + \ket{1}}{\sqrt{2}} \otimes \ket{+} = \ket{+}\ket{+}
    \end{equation*}
    Therefore, the control qubit remains in \ket{+} after the operation, confirming that the phase kickback is $0$.
    \begin{center}
        \begin{quantikz}
            \lstick{\ket{+}} & \ctrl{1} & \rstick{\ket{+}} \\
            \lstick{\ket{+}} & \gate{X} & \rstick{\ket{+}}
        \end{quantikz}
    \end{center}
    Since there is no phase kickback, the control qubit remains unchanged after the operation.
\end{examplebox}

\highspace
\begin{examplebox}[: Target in \ket{-}]
    A $CX$ gate with the target in \ket{-} will kick back a phase of $\pi$ to the control qubit. This means that the control qubit will acquire a phase of $e^{i\pi} = -1$ after the operation. Let's see if this is true:
    \begin{itemize}
        \item Control in \ket{+} to allow for superposition of control states:
        \begin{equation*}
            \ket{+} = \frac{\ket{0} + \ket{1}}{\sqrt{2}}
        \end{equation*}
        \item Target in \ket{-} to be an eigenstate of $X$ with eigenvalue $-1$:
        \begin{equation*}
            \ket{-} = \frac{\ket{0} - \ket{1}}{\sqrt{2}}
        \end{equation*}
    \end{itemize}
    The initial state is:
    \begin{equation*}
        \ket{\psi} = \ket{+}\ket{-} = \frac{\ket{0} + \ket{1}}{\sqrt{2}} \otimes \ket{-}
    \end{equation*}
    Applying the $CX$ gate (CNOT):
    \begin{center}
        \begin{quantikz}
            \lstick{\ket{+}} & \ctrl{1} & \rstick{?} \\
            \lstick{\ket{-}} & \gate{X} & \rstick{?}
        \end{quantikz}
    \end{center}
    \begin{equation*}
        CX\ket{\psi} = CX\left(\frac{\ket{0} + \ket{1}}{\sqrt{2}} \otimes \ket{-}\right) = \frac{\ket{0} \ket{-} + \ket{1} X \ket{-}}{\sqrt{2}}
    \end{equation*}
    Since $X\ket{-} = -\ket{-}$, we have:
    \begin{equation*}
        CX\ket{\psi} = \frac{\ket{0} \ket{-} + \ket{1} (-\ket{-})}{\sqrt{2}} = \frac{\ket{0} - \ket{1}}{\sqrt{2}} \otimes \ket{-} = \ket{-}\ket{-}
    \end{equation*}
    Therefore, the control qubit acquires a phase of $-1$ after the operation, confirming that the phase kickback is $\pi$.
    \begin{center}
        \begin{quantikz}
            \lstick{\ket{+}} & \ctrl{1} & \rstick{\ket{-}} \\
            \lstick{\ket{-}} & \gate{X} & \rstick{\ket{-}}
        \end{quantikz}
    \end{center}
    Since there is a phase kickback of $\pi$, the control qubit changes from \ket{+} to \ket{-} after the operation.
\end{examplebox}